\documentclass[14pt]{extreport}
\usepackage{indentfirst}
\usepackage[T2A]{fontenc}
\usepackage[utf8]{inputenc}
\usepackage[english,russian]{babel}
\usepackage{parskip}
\usepackage[top = 20mm, left = 20mm, right = 20mm, bottom = 20mm]{geometry}
\setlength\parindent{12.5mm} 
\renewcommand{\baselinestretch}{1.5} 
\def\hrf#1{\hbox to#1{\hrulefill}} 
\usepackage{ragged2e} 
\usepackage{caption}
\usepackage{tikz}
\usepackage{circuitikz}
\usetikzlibrary{decorations.pathmorphing,arrows.meta}
\usepackage{multicol}
\usepackage{float}
\usepackage{tipa}
\usepackage{graphicx}
\usepackage{enumitem}
\usepackage{amsmath}
\usepackage{amssymb}
\usepackage{listings}
\usepackage{moreverb}
\usepackage{minted}
\setcounter{tocdepth}{3}
\setcounter{secnumdepth}{3}

\lstset{
    language=C++,
    basicstyle=\ttfamily\footnotesize,
    numbers=left,               % Нумерация строк слева
    numberstyle=\tiny\color{gray}, % Стиль номеров строк
    stepnumber=1,               % Номер на каждой строке
    numbersep=5pt,              % Расстояние от номеров до кода
    breaklines=true,            % Перенос длинных строк
    tabsize=2,                  % Размер табуляции
    extendedchars=false,
    keepspaces=true,
    showstringspaces=false,
    frame=single,               % Рамка вокруг листинга
    captionpos=t,               % Подпись сверху
    xleftmargin=10pt,
    xrightmargin=5pt
}
\usepackage[colorlinks=true, citecolor=blue, urlcolor=blue, linkcolor=blue]{hyperref}
\renewcommand\thesection{\arabic{section}}


\begin{document}

\thispagestyle{empty}
    \centering{
    \hspace{0.2cm}МИНИСТЕРСТВО НАУКИ И ВЫСШЕГО ОБРАЗОВАНИЯ РОССИЙСКОЙ ФЕДЕРАЦИИ\\
    Федеральное государственное автономное образовательное учреждение высшего образования\\
    {\bf «Санкт-Петербургский политехнический университет Петра Великого»}\\[10pt]
    Институт компьютерных наук и кибербезопасности\\
    Высшая школа технологий искусственного интелекта\\
    02.03.01 Математика и компьютерные науки\\
    \vspace{1.5cm}
    {\large
    \bf Отчет по курсовой работе\\
    \bf Структуры данных}\\
    {\large\bf Разработка игрового бота для расширенной версии игры “Крестики-нолики”}\\
    \vspace{\fill}
    }
    \vspace{1.5cm}
    \flushleft{
        {\bf Выполнил:}\\
        студент группы 5130201/50301 \hspace{2.25cm} 
        \line(1,0){120} \quad Ташлыкова Е.В.\\
        {\bf Принял:}\\
        Ассистент \hspace{66mm}
        \line(1,0){120} \quad Кирпиченко С.Р.\\
        \vspace{1cm}
        \hspace{11.6cm}<<\hrf{2em}>> \hrf{6em} 2026г.
        }
        
    \vspace{1.5cm}
    \centering{
    Санкт-Петербург, 2026
    }

\newpage
\justifying
\newpage
\renewcommand\contentsname{\vspace{-3.2cm}\centering\LargeСодержание}
\tableofcontents
\thispagestyle{empty}
\newpage
{\centering\section*{Введение}}\addcontentsline{toc}{section}{Введение}
Игра «Крестики-нолики» в классическом варианте на поле 3х3 является полностью решённой игрой — при правильной игре обеих сторон партия всегда заканчивается ничьей. Однако при увеличении размера поля до 20х20, добавлении препятствий и удлинении победной линии до 5 задача выбора оптимального хода становится комбинаторно сложной. Полный перебор всех возможных продолжений невозможен из-за экспоненциального роста дерева игры.

В данной работе представлен процесс создания игрового бота для расширенной версии игры "Крестики-нолики". Она предствляет собой игровое поле размером MxN клеток, в нашем случае 20х20, с препятствиями. Выигрышной считается линия, сотстоящая из 5 и более идущих подряд символов своего игрока по вертикали, горизонтали или диагонали (под углом $45^{\circ}$).
Из-за больших размеров поля (всего 400 клеток) простой перебор для поиска подходящего хода не целесообразен, так как будет тратить слишком много времени.

Цель работы -- создание собственного игрового персонажа MyPlayer, который будет оптимально выбирать допустимую клетку для хода. Для достижения этой цели в работе формализована модель поля, описано множество рассматриваемых ходов, введены оценки клеток поля, реализована обработка немедленных побед и блокировок, а также проведены серии автоматических игр против базовых игроков. Исходный код был написан на C++.

\newpage
{\centering\section{Постановка задачи}}
Готовая игровая среда для расширенных <<Крестиков-ноликов>> уже включает в себя движок, смену ходов, историю партии и определение исхода. Всё это переписывать не требуется. Наш бот получает от среды текущее состояние игры через специальный интерфейс, а его единственная обязанность -- выдать один корректный следующий ход. Значит, нам необходимо реализовать не всю игру целиком, а только алгоритм, который выбирает, какую клетку занять текущему игроку.\\
Для того, чтобы считать работу выполненной необходимо закрыть следующие задачи:
\begin{enumerate}
    \item Разработать алгоритм для игры, удовлетворяющий критериям:
    \begin{itemize}
        \item Алгоритм должен быть эффективным, время вычисления одного хода должно занимать время порядка 50-100 мс на рабочей станции (ноутбуке);
        \item Процент побед над игроком базового уровня сложности без учета ничьих должен быть больше $50\%$; 
        \item Процент побед над игроком повышенного уровня сложности без учета ничьих должен быть больше $25\%$.
    \end{itemize}
    \item В рамках предоставленной заготовки создать своего игрока MyPlayer на языке С++;
    \item Подготовить собственные CI тесты для проверки корректности работы бота;
    \item Описать проделанную работу в итоговом отчете.
\end{enumerate}
Помимо вышеперечисленных задач важной частью работы являтеся разбиение ее этапов на конкретные issues. Для фиксирования хода выполнения курсовой работы необходимо использовать трекер задач GitHub Projects.

\newpage
{\centering\section{Математическое описание}}
{\centering\subsection{Игровое поле и допустимые ходы}}

Игровое представляет собой матрицу A размером MxN клеток. В нашем случае M=N=20. 
$(i, j)$ -- координаты клетки, при этом $ i=1,2,\ldots,N, \quad j=1,2,\ldots,M $. 

Обозначим состояние поля как функцию S(i, j), которая для каждой клетки возвращает одно из четырёх значений.
\[ S = \{  \varnothing, X, O, \# \} \]
\noindent {\Large$\varnothing$} - свободная клетка; \\
X и O -- символы игроков; \\
\# -- препятствие. 

Играют два игрока: первый (ставит X) и второй (ставит O). Игроки ходят по очереди, ставя свой знак в любую пустую клетку.
Ход допустим тогда и только тогда, когда : 
\[ 0 \leq i < N \land 0 \leq j < M \land S(i, j) = \varnothing \]

\subsection{Условие победы}
Для проверки линий на <<выигрыш>> используется четыре направления (вправо, вниз, вправо вниз по диагонали, вправо вверх по диагонали):
\[ D = \{(1,0), (0,1), (1,1), (1,-1)\} \]

Победа наступает, когда игрок собирает непрерывную линию из 5 (или более) своих символов по горизонтали, вертикали или диагонали (под углом 45°). 
Линия считается непрерывной, если между символами нет пустых клеток, чужих символов или препятствий.
Множество клеток l образует непрерывную линию в направлении $\vec{d}$, если:
\[ l = \{ (i+t*d_i, j+t*d_j) | t = 0, 1, \dots, L-1\} \]
где L - длина выигрышной линии (L=5), $\vec{d} \in D$

Для О победа наступает при наличии выигрышной линии $l_O$. Для игрока Х победа наступает, только в случае, если ответным ходом у игрока О нет $l_O$.
\[ Win_O \Longleftrightarrow \hspace{3mm} \exists \hspace{2mm} l_O \]
\[ Win_X \Longleftrightarrow \hspace{3mm} \exists \hspace{2mm} l_X \land \nexists \hspace{2mm} l_O \]
Если после хода Х на поле построенна выигрышная линия, но при этом есть открытая линия длины (L-1) символов О, то мы не можем говорить о гарантированной победе Х, так как игрок О ответным ходом может добиться ничьей.

\subsection{Кластер и его центр}
Боту важно знать, где на поле есть свободное пространство.
Кластер -- это множество пустых клеток, которые соединены друг с другом (по горизонтали, вертикали или диагонали). То есть из одной пустой клетки можно <<дойти>> до другой, шагая по соседним клеткам.
\[ C = \{ (i_1, j_1), \dots , (i_n, j_n) \}\]
Тогда центром кластера является клетка, ближайшая к геометрическому центру всех клеток множества С.
\[ c_0 = ( \frac{1}{n} \sum_{t=1}^{n} i_t \hspace{2mm};\hspace{2mm} \frac{1}{n} \sum_{t=1}^{n} j_t) \]

\subsection{Табличная оценка шаблонов}
Вместо вычисления длины линии каждый раз заново, используется предвычисленная таблица. Любая непрерывная последовательность из 5 клеток может находиться в одном из 
$ 3^5=243 $ состояний, где каждая клетка кодируется: \\
0 — пусто,
1 — свой символ,
2 — чужой символ или стена.
Для каждого из 243 шаблонов заранее вычислена числовая ценность V(шаблон) по следующей таблице \ref{tab1}:
\begin{table}[h]
\centering
\renewcommand{\tablename}{\small{Таблица.}}
\caption{\small Оценка шаблонов}
\label{tab1}
\begin{tabular}{|p{8cm}|p{5cm}|}
    \hline
    Условие & Ценность \\
    \hline
    5 своих символов & $10^8$ \\
    \hline
    4 своих символа + 1 пустая клетка & $5*10^5$ \\
    \hline
    4 своих символа + 1 непустая клетка & $5*10^4$ \\
    \hline
    3 своих символа + 2 пустых клетки по краям & $2*10^4$ \\
    \hline
    3 своих символа + 2 пустых клетки не по краям & $5*10^3$ \\
    \hline
    2 своих символа + 3 пустых клетки по краям & $5*10^2$ \\
    \hline
    2 своих символа + 3 пустыз клетки не по краям & $10^2$ \\
    \hline
    1 свой символ + 4 пустых клетки & $30$ \\
    \hline
    есть чужой символ или препятствие & $0$ \\
    \hline
\end{tabular}
\end{table}

Для клетки (i, j) и игрока s строится 4 линии по 9 клеток (4 в одну сторону, сама клетка, 4 в другую). Каждая линия разбивается на 5 скользящих окон длины L=5. Каждое окно кодируется по следующему принципу:
\[ window(l)_t = \begin{cases}
0, & A_{i_t, j_t} = \varnothing \\
1, & A_{i_t, j_t} = s \\
2, & A_{i_t, j_t} \in \{ \bar{s}, \# \} \\
\end{cases} \]
Затем окно преобразуется в индекс:
\[ index(l) = \sum_{t=0}^{4}window(l)_t \cdot 3^t \]
Далее определим $U_s(i, j)$ -- ценность клетки для игрока s.
Для каждого направления \( \vec{d} \in D \) построим расширенную линию из 9 клеток:
\[ l_9(i, j, \vec{d}) = \left\{ (i + t \cdot d_x, \; j + t \cdot d_y) \mid t = -4, -3, \dots, 0, \dots, 4 \right\}. \]
Разобьём её на 5 скользящих окон длины 5:
\[ l_9^{(k)}(i, j, \vec{d}) = \left\{ (i + (t + k)d_x, \; j + (t + k)d_y) \mid t = 0, \dots, 4 \right\}, \quad k = 0, \dots, 4. \]
Тогда:
\[ U_s(i, j) = \sum_{\vec{d} \in D} \sum_{k=0}^{4} V\bigl( \text{index}(l_9^{(k)}(i, j, \vec{d})) \bigr). \]

% \[ U_s(i, j) = \sum_{\vec{d} \in D} \sum_{n=0}^{4}V(window_{d,n})\]
% Где $window_{d,n}$ -- окно из 5 клеток, начинающееся со сдвигом n вдоль направления $\vec{d}$.
Если в двух или более направлениях обнаружены угрозы с ценностью $\geq$ 5000, итоговая сумма умножается на 10 (поощрение вилки).

\subsection{Множество ходов-кандидатов}
Для ускорения работы рассматриваются не все пустые клетки, а только перспективные. Клетка считается перспективной, если в её окрестности радиуса 2 (квадрат 5х5) есть хотя бы один символ X или O. Множество кандидатов:
\[ K = \{(i, j) = \varnothing, \exists (p, q) : |i - p| \leq 2, |j - q| \leq 2 \in \{ X, O \} \} \]

\subsection{Выбор первого хода}
Когда поле полностью пусто (на нём нет ни одного символа X или O), критерий <<близость к уже поставленным символам>> не работает. Поэтому для первого хода используется отдельная эвристика, основанная на трёх факторах:
\begin{itemize}
    \item Потенциальные линии через клетку -- количество отрезков длины L=5, которые проходят через данную клетку и целиком состоят из свободных клеток. Чем больше таких отрезков, тем больше у клетки места для построения линии в будущем.
    \item Штраф за близость к препятствиям -- если рядом с клеткой есть препятствия $(\#)$, они ограничивают возможности для роста линии. Чем ближе препятствие, тем больше штраф.
    \item Расстояние до центра крупнейшего свободного кластера -- стимулирует первый ход в сторону самой большой области пустых клеток.
\end{itemize}
Обозначим за $F_0(i, j)$ функцию оценки каждой свободной клетки при первом ходе.
\[ F_0(i, j) = 1000 L(i, j) - 20 P(i, j) - R(i, j)\]
L(i, j) -- количество отрезков длины 5, полностью проходящих через клетку (i, j) и не содержащих препятствий;\\
\[ L(i, j) = | \{ \vec{d} \in D | \forall t = 0, \dots , 4 : A_{i+td_x, j+td_y = \varnothing }\} |\]
P(i, j) -- штраф за препятствия в окрестности радиуса 2;\\
$R(i, j) = |i-i_{cluster}| + |j - j_{cluster}| $ -- манхэттенское расстояние до центра крупнейшего свободного кластера.\\

\subsection{Оценка хода}
Для каждой клетки-кандидата вычисляется итоговая оценка:
\[ F_s(i, j) = \alpha_s U_s(i, j) + \beta_{\overline{s}} U_{\overline{s}}(i, j) + B_{center}(i, j) - P(i, j) + G(i, j) \]
$s$ -- текущий игрок, $\overline{s}$ -- противник;\\
$\alpha_X = 12$ , $\beta_X = 24$ -- коэффициенты нападения и защиты (для Х);\\
$\alpha_O = 22$ , $\beta_O = 24$ -- коэффициенты нападения и защиты (для O);\\
$B_{center}(i, j)$ -- бонус центра (первые 4 хода);\\
$P(i, j)$ -- штраф за близость к препятствиям;\\
$G(i, j)$ -- бонус за близость к центру крупнейшего свободного кластера.

Вспомогательные функции:
\[ P(i, j) = \sum_{|dx| \leq 2, |dy| \leq 2, (dx, dy) \neq (0, 0)}w(|dx| + |dy|) \cdot ((i+dx, j+dy) = \#),\]
где w(1)=15, а w(2)=10 -- веса штрафа.
\[ B_{center}(i, j) = max(0, 4 - (|i-i_c| + |j-j_c|)),\]
где $(i_c, j_c)$ -- координаты центра поля (10, 10), при N=M=20.

\subsection{Итоговый алгоритм выбора хода}
На вход алгоритму подается текущее состояние поля, его размеры N и M, знак текущего игрока s. На выходе мы получаем координаты наиболее перспективного хода $(i^*, j^*)$. Последовательный план выбора подходящей клетки:
\begin{enumerate}
    \item Найти крупнейший свободный кластер (обходом в ширину).
    \item Если поле пусто, то выбрать первый ход по максимальному числу пересекающихся линий.
    \item Собрать множество кандидатов K.
    \item Если есть клетка, дающая победу (с учётом правил для X и O), то вернуть её.
    \item Если есть клетка, блокирующая победу противника, то вернуть её.
    \item Для каждой клетки из K вычислить $F_s(i, j)$
    \item Вернуть клетку с максимальной оценкой. Необходимо выбрать ходом ту клетку с координатами $(i^*, j^*)$ из множества допустимых ходов K, для которой оценочная функция $F_s$ принимает наибольшее значение
        \[(i^*, j^*) = \arg \max_{(i,j) \in K} F_s(i,j)\]
    \item Если игра за X и все оценки отрицательны, но есть ход в ничью -- вернуть его.
\end{enumerate}

\newpage
\section{Особенности реализации}
Реализация выполнена на языке C++ в классе MyPlayer. Класс содержит два поля: $m\_sign$, в котором хранится знак текущего игрока, и $m\_name$, в котором хранится имя экземпляра игрока. Метод $set\_sign$ получает знак, назначенный игровым движком, метод $get\_name$ возвращает имя игрока, а $make\_move$ принимает состояние поля и возвращает координаты следующего хода. Интерфейс класса приведён в приложении A, а полный код реализации приведён в приложении B.

Для ускорения работы используется структура FastBoard, которая хранит состояние поля в виде статического двумерного массива 20х20 элементов типа Sign. Это позволяет получать доступ к любой клетке за O(1) без накладных расходов на битовые операции исходной реализации.

Для хранения оценок используется тип long long (64-битное целое), что достаточно для значений вплоть до $10^8$.

Таблица $s\_patternScore$ размером 243 элемента является статическим членом класса. Она заполняется один раз при первом вызове $initTables()$ (вызывается в начале $make\_move$). Заполнение происходит по правилам, описанным в подразделе 2.4.
Функция $windowToIndex()$ преобразует массив из 5 чисел $\{0,1,2\}$ в индекс от 0 до 242 с использованием троичной системы счисления.

Поиск крупнейшего кластера реализован с помощью обхода в ширину (BFS) по 8 направлениям (включая диагонали). Для каждой компоненты:
\begin{itemize}
    \item Подсчитывается размер.
    \item Вычисляется геометрический центр как среднее арифметическое координат.
    \item Выбирается клетка компоненты, ближайшая к этому центру (по евклидову расстоянию). При равенстве учитывается количество возможных линий длины 5 через клетку и штраф за препятствия.
\end{itemize}

Для повышения качества принимаемых решений в алгоритме используется ограниченный поиск в глубину на основе негамакса с альфа-бета отсечением. Негамакс — это компактная форма минимакса, в которой оценка позиции для текущего игрока вычисляется как максимум из отрицательных оценок позиций после каждого возможного хода.
Для оптимизации алгоритма $negamax$ с одноименной функцией его реализации в коде создана функция $getDynamicDepth$, которая определяет глубину ходов, если ценность клетки  $U_s(i, j) \geq 2.000.000$, то глубина увеличивается до 5, если $U_s(i, j) \geq 50.000$, то до 4, иначе -- глубина 3.

Основной цикл $make\_move$:
\begin{minted}[fontsize=\footnotesize]{c}
Point MyPlayer::make_move(const State &state){
        initTables();
        FastBoard board;
        board.sync(state);
        int moveNumber = state.get_move_no();
        ClusterInfo cluster = findLargestCluster(board);
        std::vector<Point> cells = getCandidateCells(board);
        if (cells.empty())
          return {0, 0}; 
        if (moveNumber == 0)
          return chooseFirstMove(board, cluster);
        Sign opponent = (m_sign == Sign::X) ? Sign::O : Sign::X;
        // немедленная победа 
        if (m_sign == Sign::O){
          for (const auto &cell : cells){
            if (hasLineAfterMove(board, cell.x, cell.y, Sign::O))
              return cell;
          }
        }
        else{
          for (const auto &cell : cells){
            if (isRealXWin(board, cell.x, cell.y))
              return cell;
          }
        }
        // защита от победы противника
        for (const auto &cell : cells){
          if (hasLineAfterMove(board, cell.x, cell.y, opponent))
            return cell;
        }
        // запоминаем ничейный ход для X 
        Point drawMove = cells[0];
        bool hasDraw = false;
        if (m_sign == Sign::X){
          for (const auto &cell : cells){
            if (isXDraw(board, cell.x, cell.y)){
              drawMove = cell;
              hasDraw = true;
              break;
            }
          }
        }
        // жадная оценка и отбор кандидатов
        std::vector<RatedMove> rootMoves;
        for (const auto &cell : cells){
          long long score = evaluateCell(board, cell.x, cell.y, cluster, moveNumber);
          // проверка на мгновенную победу 
          if (valueScore(board, m_sign, cell.x, cell.y) >= WIN_SCORE)
            return cell;
          rootMoves.push_back({cell.x, cell.y, score});
        }
        std::sort(rootMoves.begin(), rootMoves.end(), [](const RatedMove &a, const RatedMove &b)
                  { return a.weight > b.weight; });
        if (rootMoves.size() > static_cast<size_t>(SEARCH_WIDTH))
          rootMoves.resize(SEARCH_WIDTH);
        // негамакс на кандидатах 
        int depth = getDynamicDepth(board, m_sign);
        long long bestScore = -WIN_SCORE * 2;
        Point bestMove = {rootMoves[0].x, rootMoves[0].y};
        for (const auto &move : rootMoves){
          Sign oldValue = board.get(move.x, move.y);
          board.set(move.x, move.y, m_sign);
          long long score = -negamax(board, depth - 1, -WIN_SCORE * 2, WIN_SCORE * 2,
                                     opponent, move.x, move.y, moveNumber + 1);
          board.set(move.x, move.y, oldValue);
          if (score > bestScore){
            bestScore = score;
            bestMove = {move.x, move.y}; 
          }
        }
        if (m_sign == Sign::X && hasDraw && bestScore < 0) // уход в ничью для X 
          return drawMove;
        return bestMove;
      }
\end{minted}


\newpage
\section{Результаты работы программы}
Проверка бота проходила в рамках тестирования в стандартном игровом движке проекта на поле 20х20 с исходными правилами победы и препятствий. Было сыграно по 500 партий, где MyPlayer играл за X, против данного изначально базового игрока простого и сложного уровней. Дополнительно прошла серия из 500 игр MyPlayer против самого себя.
В таблице \ref{tab2} приведены абсолютные и процентные результаты игр. Проценты рассчитаны относительно 500 партий в каждой серии.
\begin{table}[h]
\centering
\renewcommand{\tablename}{\small{Таблица.}}
\caption{\small Результаты серийного тестирования бота MyPlayer}
\label{tab2}
\begin{footnotesize}
    \begin{tabular}{|p{5cm}|p{2cm}|p{2cm}|p{2cm}|p{2cm}|}
    \hline
    Серия & Победы MyPlayer & Победы соперника & Ничьи & Победы без ничьих \\
    \hline
    MyPlayer -- BaselineEasy & $429 / 85.8 \%$ & $43 / 8.6 \%$ & $28 / 5.6 \%$ & $90.89 \%$ \\
    \hline
    MyPlayer -- BaselineHard & $112 / 22.4 \%$ & $73 / 14.6 \%$ & $315 / 63 \%$ & $60.54 \%$ \\
    \hline
    MyPlayer -- MyPlayer & $195 / 39 \%$ & $95 / 19 \%$ & $210 / 42 \%$ & $67.24 \%$ \\
    \hline
    \end{tabular}
\end{footnotesize}
\end{table}

Время выбора хода измерялось средствами игрового движка. Работа была выполнена на ноутбуке MacBook Pro M4. Средние значения приведены в таблице \ref{tab3}.
\begin{table}[h]
\centering
\renewcommand{\tablename}{\small{Таблица.}}
\caption{\small Среднее время расчёта хода и обработки партии}
\label{tab3}
\begin{footnotesize}
    \begin{tabular}{|p{5cm}|p{3cm}|p{3cm}|p{3cm}|}
    \hline
    Серия & Ход MyPlayer (мс) & Ход соперника (мс) & Средняя продолжительность игры (мс) \\
    \hline
    MyPlayer -- BaselineEasy & $82.3257$ & $0.0006$ & $41.1668$ \\
    \hline
    MyPlayer -- BaselineHard & $66.7691$ & $0.0007$ & $33.3877$ \\
    \hline
    MyPlayer -- MyPlayer & $66.9879$ & $65.8646$ & $66.4288$ \\
    \hline
    \end{tabular}
\end{footnotesize}
\end{table}

\newpage
{\centering\section*{Заключение}}\addcontentsline{toc}{section}{Заключение}
В результате выполнения курсовой работы был разработан и реализован игровой бот для усложненных <<Крестиков-ноликов>>. В основе алгоритма лежит комбинация метода negamax с alpha-beta отсечением и эвристической функцией оценки состояния доски. Для повышения производительности применялись низкоуровневые оптимизации: битовое представление игрового поля и предвычисление таблицы линий на этапе компиляции.

Проведённое тестирование показало, что разработанный алгоритм является достаточно сильным. Бот уверенно превосходит простого соперника (результат побед без учета ничьих составляет $90.89 \%$). При игре против сложного бота результаты оказались также выше необходимых -- $60.54 \%$ побед без учета ничьих.

Также алгоритм можно назвать достаточно эффективным, так как время вычисления одного хода занимает менее 100мс, а именно -- 82.33мс при игре с BaselineEasy и 66.77мс при игре с BaselineHard.

В процессе работы были детально изучены базовые алгоритмы теории игр: minimax, negamax, alpha-beta pruning, а также практические подходы к оптимизации переборных алгоритмов. Полученные навыки включают работу с матричными играми и проектированием собственных эвристических функций.

\newpage
\renewcommand{\bibname}{\vspace{-3.2cm}\centering\LargeСписок используемых источников}\addcontentsline{toc}{section}{Список используемых источников}
\begin{thebibliography}{8}

\bibitem{gametheory}
Петросян Л. А., Зенкевич Н. А., Семина Е. А. Теория игр : учеб. пособие для ун-тов. -- Москва : Высш. шк. : Книжный дом «Университет», 1998. - 304 с. - ISBN 5-06-001005-8.

\bibitem{sennov_tictactoe}
Сеннов В. tictactoe-course, материалы к курсовой работе [Электронный ресурс] : 
\href{https://github.com/vsennov/tictactoe-course}{https://github.com/vsennov/tictactoe-course} 
(дата обращения: 22.05.2026).

\bibitem{habr_chess}
Habr. <<Как компьютер играет в шахматы?>> [Электронный ресурс]: 
\href{https://habr.com/ru/articles/390821/}{https://habr.com/ru/articles/390821/} 
(дата обращения: 20.05.2026).

\bibitem{habr_optimization}
Habr. <<Оптимизация перебора>> [Электронный ресурс] : 
\href{https://habr.com/ru/articles/190850/}{https://habr.com/ru/articles/190850/} 
(дата обращения: 22.05.2026).

\bibitem{devto_projects}
Adiati, A. <<Planning and Tracking Projects with GitHub's Projects Tool>> [Электронный ресурс] : 
\href{https://dev.to/adiatiayu/planning-and-tracking-projects-with-githubs-projects-tool-572l}{https://dev.to/adiatiayu/planning-and-tracking-projects-with-githubs-projects-tool-572l} 
(дата обращения: 23.05.2026).

\end{thebibliography}

\newpage
\captionsetup[listing]{labelformat=empty}
{\flushleft\section*{Приложение А. Исходный код заголовочного файла $my\_player.hpp$}}\addcontentsline{toc}{section}{A. Исходный код заголовочного файла $my\_player.hpp$}
\lstinputlisting[
    caption={Файл my\_player.hpp},
    label=lst:my_player_hpp
]{my_player.hpp}

\newpage
\captionsetup[listing]{labelformat=empty}
{\flushleft\section*{Приложение В. Исходный код файла реализации $my\_player.cpp$}}\addcontentsline{toc}{section}{B. Исходный код файла реализации $my\_player.cpp$}
\lstinputlisting[
    caption={Файл my\_player.cpp},
    label=lst:my_player_cpp
]{my_player.cpp}

\newpage
{\flushleft\section*{Приложение C. Трекинг задач, Kanban и сопровождение курсовой работы}}\addcontentsline{toc}{section}{C. Трекинг задач, Kanban и сопровождение курсовой работы}
\begin{figure}[h!]
  \centering
  \includegraphics[height=10cm]{project.png}
  \caption{\small Снимок Kanban-доски}
\end{figure}
\begin{figure}[h!]
  \centering
  \includegraphics[height=15cm]{git.png}
  \caption{\small Снимок графа коммитов}
\end{figure}

\end{document}
